\section*{Description}
Ce document est né de mon envie de rédiger un cours de mathématiques à ma façon. Ainsi, il se doit d'être long, d'investiguer chaque hypothèse et de tenter de donner l'intuition derrière certaines abstractions mathématiques. On démontrera tout les résultats énoncés, même si la démonstration prendra parfois beaucoup de temps, car il faudrait développer beaucoup de théorie. \\

Malgré le fait que ce document ne soit pas un "vrai" document de cours, j'espère (et cela peut se sentir dans la rédaction) que ce document sera utilisé par des personnes qui souhaiterait utiliser mon approche des mathématiques. Cependant, il est évident que ce document n'est pas fait pour une première approche en la matière. Il est écrit après avoir suivi des cours, fait beaucoup d'exercices mais surtout après avoir consommé énormément de contenus mathématiques. \\

Pour le fonctionnement de ce document, il fonctionne comme un algorithme dit de "parcours en profondeur". On fera toujours le choix d'explorer tout les horizons raisonnables donnés par la théorie avant de faire un bon en avant. Je le publie sur \href{https://github.com/pierremechelinck/maths}{\texttt{github.com}} pour pouvoir le consulter rapidement et le partager. \\


\emph{J’autorise quiconque le souhaite à placer sur son site un lien vers ce document, mais je n’autorise personne à l' héberger directement. J’interdis par ailleurs toute utilisation commerciale de mon travail.}

\section*{Motivations}
La première motivation de ce texte était de réviser mes fondements et d'écrire un cours de mathématiques pour moi, à ma façon. J'avais vraiment envie d'avoir un document qui couvrait les cours que j'ai suivi jusqu'à présent. \\

La deuxième motivation, venue plus tard, est fondée sur un avis personnel. Je suis de l'avis que l'enseignement des mathématiques fondamentales, pour les individus qui veulent les poursuivre, est retenu par son utilité dans les sciences appliquées. C'est ainsi que je fais un cours de mathématiques pour rien; au sens où j'aime les maths en tant que tel.


\epigraph{We must know — we will know!}{\textit{David Hilbert (1862-1943)}}


\section*{Remerciements}
J'aimerais remercier les enseignants de l'université de Rouen Normandie pour les cours que j'ai pu y suivre en première année de licence. J'aimerais aussi remercier les enseignants de l'Ecole universitaire de premier cycle de Paris Saclay pour les cours.  \\

Ce cours a été en partie rédigé à partir de notes de cours et de documents disponibles librement sur Internet. Je remercie ainsi les individus et les institutions qui laissent leurs cours disponibles gratuitement en ligne. \\


Pour finir, j'aimerais remercier d'avance les personnes qui m'enverront des corrections pour les énoncés éventuellement erronés de ce polycopié. 

















